\section{Introduction}\label{sec:intro}
Most contemporary machine learning systems are tailored to a specific data modality: convolutional networks for images, recurrent or transformer networks for sequences, graph neural networks for relational data.
Supervised learning dominates, requiring large amounts of labelled data that embed human conceptual biases.

The APEIRON Framework~\cite{Beniers2026Apeiron} takes a radically different approach: it builds knowledge from elementary, unlabelled observables through a sequence of emergent layers.
In earlier work, the Layer~2 relational hypergraph and Layer~3 functional clustering pipeline was validated on benchmark vision datasets (MNIST, Fashion‑MNIST, Kuzushiji‑MNIST, CIFAR‑10), where it autonomously discovered spatially and semantically coherent pixel groups~\cite{Beniers2026Empirical}.

The core question of the present paper is: **does this pipeline work only on images, or is it genuinely domain‑independent?**
To answer this, we apply the identical Layer~2--3 pipeline to three fundamentally different types of time‑series data:
\begin{itemize}
\item \textbf{Electrocardiograms (ECG)} – physiological signals with known cyclic structure.
\item \textbf{Financial time series} – daily stock returns, where hidden sector correlations are often hypothesised.
\item \textbf{Seismic records} – continuous ground‑motion signals exhibiting event‑related patterns.
\end{itemize}
In each case, we construct a co‑activation graph from short sliding‑window segments, apply Louvain community detection, and evaluate the functional relevance of the discovered clusters via counterfactual ablation.

The results demonstrate that the pipeline discovers modules that correspond to physiologically meaningful ECG components (QRS complex), financial sectors, and seismic event phases – without any domain knowledge or labels.
This provides strong evidence that the APEIRON substrate is a universal, modality‑agnostic structure‑discovery engine.